%%%%%%%%%%%%%%%%%%%%%%%%%%%%%%%%%%%%%%%%%%%%%%%%%%%%%%%%%%%%%%%%%%%%%%
% Lớp 8 | Chương III. Tứ giác | Hình chữ nhật
% Dạng 8H3?13-3: Tính độ dài, bài toán thực tế
%%%%%%%%%%%%%%%%%%%%%%%%%%%%%%%%%%%%%%%%%%%%%%%%%%%%%%%%%%%%%%%%%%%%%%


% nguồn: De_AI_Toan8_cuoi_hoc_ki_i_1 (nhập 21/11/2026)
\begin{bt}%[8H3H13-3]%[ma:5dc2840c]
	Một sân thể thao hình chữ nhật có chiều dài $16\text{ m}$ và chiều rộng $12\text{ m}$.
a) Tính chu vi và diện tích của sân thể thao.
b) Người ta muốn căng một sợi dây nối hai góc đối diện theo đường chéo của sân. Tính chiều dài tối thiểu của sợi dây đó.
	\loigiai{
		a) Chu vi sân thể thao là $2 \cdot (16 + 12) = 56\text{ m}$. \\
Diện tích sân thể thao là $16 \cdot 12 = 192\text{ m}^2$. \\
b) Độ dài đường chéo $d$ của sân hình chữ nhật được tính theo định lí Pythagore: \\
$d^2 = 16^2 + 12^2 = 256 + 144 = 400 \Rightarrow d = 20\text{ m}$. \\
Vậy sợi dây cần có chiều dài tối thiểu là $20\text{ m}$.
	}
\end{bt}

% nguồn: De_AI_Toan8_cuoi_hoc_ki_i_1 (nhập 21/11/2026)
\begin{bt}%[8H3H13-3]%[ma:f5ca790e]
	Cho hình chữ nhật $ABCD$ có chu vi bằng $34\text{ cm}$ và độ dài đường chéo $AC = 13\text{ cm}$. Tính diện tích của hình chữ nhật $ABCD$.
	\loigiai{
		Gọi $AB = a$, $BC = b$ ($a, b > 0$). Vì chu vi hình chữ nhật là $34\text{ cm}$ nên $a + b = \dfrac{34}{2} = 17\text{ cm}$. \\
Áp dụng định lí Pythagore trong tam giác $ABC$ vuông tại $B$, ta có $a^2 + b^2 = AC^2 = 13^2 = 169$. \\
Áp dụng hằng đẳng thức $(a + b)^2 = a^2 + 2ab + b^2$, ta suy ra: \\
$2ab = (a + b)^2 - (a^2 + b^2) = 17^2 - 169 = 289 - 169 = 120$. \\
Do đó $ab = \dfrac{120}{2} = 60$. Vậy diện tích hình chữ nhật $ABCD$ là $60\text{ cm}^2$.
	}
\end{bt}
